% Curs 12 - Pagina 4 (3-I-1995).
% Fragment pentru % Curs 12 - Pagina 4 (3-I-1995).
% Fragment pentru % Curs 12 - Pagina 4 (3-I-1995).
% Fragment pentru % Curs 12 - Pagina 4 (3-I-1995).
% Fragment pentru \input{courses/pages/curs12-pagina4-ro}.
% Transcriere in ordinea manuscrisului; fara preambul si fara salt fortat de pagina.

\subsection*{Curs 12 - Pagina 4}

\noindent\underline{Dem.}\quad Fie $u\in A^+$,
$u=x_1x_2\cdots x_m$, $x_i\in A$.

\[
f(u)\ \overset{?}{=}\ \zeta(x_1)\zeta(x_2)\cdots\zeta(x_m)\in S.
\]
% Semnul intrebarii deasupra egalitatii apare in manuscris,
% inaintea verificarii ca aplicatia astfel propusa este un morfism.

\noindent $v\in A^+$, $v=y_1y_2\cdots y_n$.
\[
f(v)=\zeta(y_1)\cdots\zeta(y_n).
\]
\[
f(uv)=\zeta(x_1)\cdots\zeta(x_m)\zeta(y_1)\cdots\zeta(y_n)
=f(u)f(v).
\]

\[
\begin{tikzcd}[column sep=large,row sep=large]
A \arrow[r,"\zeta_A"] \arrow[dr,"\zeta"'] & A^+ \arrow[d,"f"] \\
& S
\end{tikzcd}
\]

\medskip
\noindent $(!)$ Lu\u{a}m $g:A^+\longrightarrow S$ morfism de semigrup
\c{s}i $g\circ\zeta_A=\zeta\Rightarrow g=f$.

\medskip
\noindent $u\in A^+$,
\[
u=x_1\cdots x_m=\zeta_A(x_1)\cdots\zeta_A(x_m).
\]
\begin{align*}
g(u)
 &=g\bigl(\zeta_A(x_1)\cdots\zeta_A(x_m)\bigr)\\
 &=g\bigl(\zeta_A(x_1)\bigr)\cdots g\bigl(\zeta_A(x_m)\bigr)\\
 &=\zeta(x_1)\cdots\zeta(x_m)=f(u),
 \qquad (\forall)\,u\in A^+.
\end{align*}

\bigskip
\noindent\fbox{T.}\quad\underline{Pentru monoidul liber}

\medskip
\noindent $A\neq\varnothing$, $A^+$,
\quad $A\xrightarrow{\ \zeta_A\ }A^+$.

\medskip
\noindent $(\forall)$ monoid $(M,\cdot,e)$, $(\forall)$ aplica\c{t}ie de mul\c{t}imi
$\zeta:A\longrightarrow M$,
\[
\begin{tikzcd}[column sep=large,row sep=large]
A \arrow[r,"\zeta_A"] \arrow[d,"\zeta"'] & A^+ \arrow[dl,dashed,"\exists!\,f"] \\
M &
\end{tikzcd}
\]
\noindent $\exists!\,f:A^+\longrightarrow M$ a.\^{\i}. $f\circ\zeta_A=\zeta$.

\medskip
\noindent\underline{Dem.}\quad Analog $+$ $f(1)\overset{\mathrm{def}}{=}e\in M$.
% NOTA EDITORIALA: Titlul spune „monoidul liber” si se defineste f(1),
% dar in enunt si in diagrama apare A^+, nu A^*. Am pastrat lectura manuscrisului;
% pentru proprietatea universala a monoidului liber s-ar astepta A^*.

\bigskip
\noindent\underline{Congruen\c{t}e pe $A^+$ generate de o rela\c{t}ie binar\u{a}}
$R\subseteq A^+\times A^+$.

\[
R=\bigl\{(w_i,w_i')\in A^+\times A^+\mid i\in I\bigr\}.
\]
% Nota: litera care indexeaza perechile apare ca I (sau un simbol apropiat).

\noindent $(1)$\quad
$R^{-1}=\bigl\{(v,u)\mid (u,v)\in R\bigr\}$.
% Pagina se incheie aici; continuarea definitiilor ramane pentru pagina urmatoare.
.
% Transcriere in ordinea manuscrisului; fara preambul si fara salt fortat de pagina.

\subsection*{Curs 12 - Pagina 4}

\noindent\underline{Dem.}\quad Fie $u\in A^+$,
$u=x_1x_2\cdots x_m$, $x_i\in A$.

\[
f(u)\ \overset{?}{=}\ \zeta(x_1)\zeta(x_2)\cdots\zeta(x_m)\in S.
\]
% Semnul intrebarii deasupra egalitatii apare in manuscris,
% inaintea verificarii ca aplicatia astfel propusa este un morfism.

\noindent $v\in A^+$, $v=y_1y_2\cdots y_n$.
\[
f(v)=\zeta(y_1)\cdots\zeta(y_n).
\]
\[
f(uv)=\zeta(x_1)\cdots\zeta(x_m)\zeta(y_1)\cdots\zeta(y_n)
=f(u)f(v).
\]

\[
\begin{tikzcd}[column sep=large,row sep=large]
A \arrow[r,"\zeta_A"] \arrow[dr,"\zeta"'] & A^+ \arrow[d,"f"] \\
& S
\end{tikzcd}
\]

\medskip
\noindent $(!)$ Lu\u{a}m $g:A^+\longrightarrow S$ morfism de semigrup
\c{s}i $g\circ\zeta_A=\zeta\Rightarrow g=f$.

\medskip
\noindent $u\in A^+$,
\[
u=x_1\cdots x_m=\zeta_A(x_1)\cdots\zeta_A(x_m).
\]
\begin{align*}
g(u)
 &=g\bigl(\zeta_A(x_1)\cdots\zeta_A(x_m)\bigr)\\
 &=g\bigl(\zeta_A(x_1)\bigr)\cdots g\bigl(\zeta_A(x_m)\bigr)\\
 &=\zeta(x_1)\cdots\zeta(x_m)=f(u),
 \qquad (\forall)\,u\in A^+.
\end{align*}

\bigskip
\noindent\fbox{T.}\quad\underline{Pentru monoidul liber}

\medskip
\noindent $A\neq\varnothing$, $A^+$,
\quad $A\xrightarrow{\ \zeta_A\ }A^+$.

\medskip
\noindent $(\forall)$ monoid $(M,\cdot,e)$, $(\forall)$ aplica\c{t}ie de mul\c{t}imi
$\zeta:A\longrightarrow M$,
\[
\begin{tikzcd}[column sep=large,row sep=large]
A \arrow[r,"\zeta_A"] \arrow[d,"\zeta"'] & A^+ \arrow[dl,dashed,"\exists!\,f"] \\
M &
\end{tikzcd}
\]
\noindent $\exists!\,f:A^+\longrightarrow M$ a.\^{\i}. $f\circ\zeta_A=\zeta$.

\medskip
\noindent\underline{Dem.}\quad Analog $+$ $f(1)\overset{\mathrm{def}}{=}e\in M$.
% NOTA EDITORIALA: Titlul spune „monoidul liber” si se defineste f(1),
% dar in enunt si in diagrama apare A^+, nu A^*. Am pastrat lectura manuscrisului;
% pentru proprietatea universala a monoidului liber s-ar astepta A^*.

\bigskip
\noindent\underline{Congruen\c{t}e pe $A^+$ generate de o rela\c{t}ie binar\u{a}}
$R\subseteq A^+\times A^+$.

\[
R=\bigl\{(w_i,w_i')\in A^+\times A^+\mid i\in I\bigr\}.
\]
% Nota: litera care indexeaza perechile apare ca I (sau un simbol apropiat).

\noindent $(1)$\quad
$R^{-1}=\bigl\{(v,u)\mid (u,v)\in R\bigr\}$.
% Pagina se incheie aici; continuarea definitiilor ramane pentru pagina urmatoare.
.
% Transcriere in ordinea manuscrisului; fara preambul si fara salt fortat de pagina.

\subsection*{Curs 12 - Pagina 4}

\noindent\underline{Dem.}\quad Fie $u\in A^+$,
$u=x_1x_2\cdots x_m$, $x_i\in A$.

\[
f(u)\ \overset{?}{=}\ \zeta(x_1)\zeta(x_2)\cdots\zeta(x_m)\in S.
\]
% Semnul intrebarii deasupra egalitatii apare in manuscris,
% inaintea verificarii ca aplicatia astfel propusa este un morfism.

\noindent $v\in A^+$, $v=y_1y_2\cdots y_n$.
\[
f(v)=\zeta(y_1)\cdots\zeta(y_n).
\]
\[
f(uv)=\zeta(x_1)\cdots\zeta(x_m)\zeta(y_1)\cdots\zeta(y_n)
=f(u)f(v).
\]

\[
\begin{tikzcd}[column sep=large,row sep=large]
A \arrow[r,"\zeta_A"] \arrow[dr,"\zeta"'] & A^+ \arrow[d,"f"] \\
& S
\end{tikzcd}
\]

\medskip
\noindent $(!)$ Lu\u{a}m $g:A^+\longrightarrow S$ morfism de semigrup
\c{s}i $g\circ\zeta_A=\zeta\Rightarrow g=f$.

\medskip
\noindent $u\in A^+$,
\[
u=x_1\cdots x_m=\zeta_A(x_1)\cdots\zeta_A(x_m).
\]
\begin{align*}
g(u)
 &=g\bigl(\zeta_A(x_1)\cdots\zeta_A(x_m)\bigr)\\
 &=g\bigl(\zeta_A(x_1)\bigr)\cdots g\bigl(\zeta_A(x_m)\bigr)\\
 &=\zeta(x_1)\cdots\zeta(x_m)=f(u),
 \qquad (\forall)\,u\in A^+.
\end{align*}

\bigskip
\noindent\fbox{T.}\quad\underline{Pentru monoidul liber}

\medskip
\noindent $A\neq\varnothing$, $A^+$,
\quad $A\xrightarrow{\ \zeta_A\ }A^+$.

\medskip
\noindent $(\forall)$ monoid $(M,\cdot,e)$, $(\forall)$ aplica\c{t}ie de mul\c{t}imi
$\zeta:A\longrightarrow M$,
\[
\begin{tikzcd}[column sep=large,row sep=large]
A \arrow[r,"\zeta_A"] \arrow[d,"\zeta"'] & A^+ \arrow[dl,dashed,"\exists!\,f"] \\
M &
\end{tikzcd}
\]
\noindent $\exists!\,f:A^+\longrightarrow M$ a.\^{\i}. $f\circ\zeta_A=\zeta$.

\medskip
\noindent\underline{Dem.}\quad Analog $+$ $f(1)\overset{\mathrm{def}}{=}e\in M$.
% NOTA EDITORIALA: Titlul spune „monoidul liber” si se defineste f(1),
% dar in enunt si in diagrama apare A^+, nu A^*. Am pastrat lectura manuscrisului;
% pentru proprietatea universala a monoidului liber s-ar astepta A^*.

\bigskip
\noindent\underline{Congruen\c{t}e pe $A^+$ generate de o rela\c{t}ie binar\u{a}}
$R\subseteq A^+\times A^+$.

\[
R=\bigl\{(w_i,w_i')\in A^+\times A^+\mid i\in I\bigr\}.
\]
% Nota: litera care indexeaza perechile apare ca I (sau un simbol apropiat).

\noindent $(1)$\quad
$R^{-1}=\bigl\{(v,u)\mid (u,v)\in R\bigr\}$.
% Pagina se incheie aici; continuarea definitiilor ramane pentru pagina urmatoare.
.
% Transcriere in ordinea manuscrisului; fara preambul si fara salt fortat de pagina.

\subsection*{Curs 12 - Pagina 4}

\noindent\underline{Dem.}\quad Fie $u\in A^+$,
$u=x_1x_2\cdots x_m$, $x_i\in A$.

\[
f(u)\ \overset{?}{=}\ \zeta(x_1)\zeta(x_2)\cdots\zeta(x_m)\in S.
\]
% Semnul intrebarii deasupra egalitatii apare in manuscris,
% inaintea verificarii ca aplicatia astfel propusa este un morfism.

\noindent $v\in A^+$, $v=y_1y_2\cdots y_n$.
\[
f(v)=\zeta(y_1)\cdots\zeta(y_n).
\]
\[
f(uv)=\zeta(x_1)\cdots\zeta(x_m)\zeta(y_1)\cdots\zeta(y_n)
=f(u)f(v).
\]

\[
\begin{tikzcd}[column sep=large,row sep=large]
A \arrow[r,"\zeta_A"] \arrow[dr,"\zeta"'] & A^+ \arrow[d,"f"] \\
& S
\end{tikzcd}
\]

\medskip
\noindent $(!)$ Lu\u{a}m $g:A^+\longrightarrow S$ morfism de semigrup
\c{s}i $g\circ\zeta_A=\zeta\Rightarrow g=f$.

\medskip
\noindent $u\in A^+$,
\[
u=x_1\cdots x_m=\zeta_A(x_1)\cdots\zeta_A(x_m).
\]
\begin{align*}
g(u)
 &=g\bigl(\zeta_A(x_1)\cdots\zeta_A(x_m)\bigr)\\
 &=g\bigl(\zeta_A(x_1)\bigr)\cdots g\bigl(\zeta_A(x_m)\bigr)\\
 &=\zeta(x_1)\cdots\zeta(x_m)=f(u),
 \qquad (\forall)\,u\in A^+.
\end{align*}

\bigskip
\noindent\fbox{T.}\quad\underline{Pentru monoidul liber}

\medskip
\noindent $A\neq\varnothing$, $A^+$,
\quad $A\xrightarrow{\ \zeta_A\ }A^+$.

\medskip
\noindent $(\forall)$ monoid $(M,\cdot,e)$, $(\forall)$ aplica\c{t}ie de mul\c{t}imi
$\zeta:A\longrightarrow M$,
\[
\begin{tikzcd}[column sep=large,row sep=large]
A \arrow[r,"\zeta_A"] \arrow[d,"\zeta"'] & A^+ \arrow[dl,dashed,"\exists!\,f"] \\
M &
\end{tikzcd}
\]
\noindent $\exists!\,f:A^+\longrightarrow M$ a.\^{\i}. $f\circ\zeta_A=\zeta$.

\medskip
\noindent\underline{Dem.}\quad Analog $+$ $f(1)\overset{\mathrm{def}}{=}e\in M$.
% NOTA EDITORIALA: Titlul spune „monoidul liber” si se defineste f(1),
% dar in enunt si in diagrama apare A^+, nu A^*. Am pastrat lectura manuscrisului;
% pentru proprietatea universala a monoidului liber s-ar astepta A^*.

\bigskip
\noindent\underline{Congruen\c{t}e pe $A^+$ generate de o rela\c{t}ie binar\u{a}}
$R\subseteq A^+\times A^+$.

\[
R=\bigl\{(w_i,w_i')\in A^+\times A^+\mid i\in I\bigr\}.
\]
% Nota: litera care indexeaza perechile apare ca I (sau un simbol apropiat).

\noindent $(1)$\quad
$R^{-1}=\bigl\{(v,u)\mid (u,v)\in R\bigr\}$.
% Pagina se incheie aici; continuarea definitiilor ramane pentru pagina urmatoare.
